\documentclass[aspectratio=169,usepdftitle=true,11pt,ngerman,t]{beamer}

\newcommand\lqfirstname{Name}
\newcommand\lqfullname{Name Last}
\newcommand\lqemail{name.last}
\errorcontextlines999999

% base
\usepackage[utf8]{inputenc}
\usepackage[T1]{fontenc}
\usepackage{lmodern}
\usepackage{microtype}
\usepackage[main=ngerman]{babel}
\usepackage{csquotes}

% graphics
\usepackage[prefix=]{xcolor-material}
\usepackage{graphics}
\usepackage{tcolorbox}
\usepackage{twemojis}
\usepackage{tikz}
\usetikzlibrary{
    positioning,
    calc,
    tikzmark,
    circuits.logic.IEC
}

% math
\usepackage{amsmath,amssymb}
\usepackage{mathtools}
\usepackage{array}
\usepackage{booktabs}
\usepackage{cancel}
\usepackage{ulem}

% flo
\usepackage{fancyqr}
\usepackage[addons]{color-palettes}
\usepackage[encoding,noloadlangs,cpalette,numinpar,fakeminted]{sopra-listings}
\usepackage{code-animation}
\usetheme[theorems]{colorful-dream}
\theoremstyle{exercise}
\newtheorem*{exercise*}{Aufgabe}
\theoremstyle{solve}
\newtheorem*{solve*}{Lösung}

% tools
\usepackage{algpseudocode}
\usepackage{dirtree}

% warnings
\usepackage{silence}
\WarningFilter{latexfont}{Some font shapes were not available}
\WarningFilter{latexfont}{Font shape}
\WarningFilter{pgf}{Returning node center instead of a point on node border}
\let\origenddocument\enddocument\renewcommand{\enddocument}{\hfuzz=10.95pt\origenddocument}
% title image: box
% ([vshift], inverse scale, width, content)
\newcommand\titleimagebox[4][0pt]{%
\newlength\tibheight\setlength\tibheight{4cm}
\newlength\tibtY\setlength\tibtY{-6em}
\titleimage{\begin{minipage}[c][#2\tibheight]{#3}\vspace*{#2\tibtY}%
    \vspace*{#1}\begin{tcolorbox}[boxrule=0.1pt]%
    \centering#4% TODO: bug: not centered if too wide
    \end{tcolorbox}%
\end{minipage}}}

% outro image: text
% (text)
\newcommand\outroimagetext[1]{\outro{Karlsruhe, \today%
\begin{tikzpicture}[remember picture,overlay]
    \node[align=center] at (current page.center) {#1};
\end{tikzpicture}}}% chktex 31

% reset counters
\newcommand\resetframecounters{\addtocounter{exercise}{-1}\addtocounter{solve}{-1}}
\newcommand\resetex{\addtocounter{exercise}{-1}}
\newcommand\resetsolve{\addtocounter{solve}{-1}}

% small footnote
\newcommand{\nicefootnote}[1]{\tikz[overlay,remember picture]{\node[anchor=south,text width=\dimexpr\textwidth+2em\relax,align=center,above=1em of current page.south] (nicefootnote) {\hrule\vspace{1pt}\footnotesize#1};}}

% kv
% ([draw], variables, values)
% {currently only works correctly for 4 variables}
\newcommand{\lqkv}[3][]{
    \begin{tikzpicture}
        \foreach \v [count=\k from 1] in {#2} {\global\let\Vdim=\k}
        \pgfmathsetmacro\rows{2^ceil(\Vdim/2)}
        \pgfmathsetmacro\cols{2^int(\Vdim/2)}

        \foreach \v [count=\k from 1] in {#2} {
            \pgfmathsetmacro\shift{0.5+floor(\k*ceil(\Vdim/2)/2)}

            \ifcase\k\relax
            \or
                \draw (\shift-1,0.75) -- (\shift+1,0.75);
                \node[fill=white] at (\shift,0.75) {\v};
            \or
                \draw (-0.75,-\shift) -- (-0.75,-\shift+2);
                \node[fill=white] at (-0.75,-\shift+1) {\v};
            \or
                \draw (\shift-2,-\cols+0.25) -- (\shift,-\cols+0.25);
                \node[fill=white] at (\shift-1,-\cols+0.25) {\v};
            \or
                \draw (\cols-0.25,-\shift+1) -- (\cols-0.25,-\shift+3);
                \node[fill=white] at (\cols-0.25,-\shift+2) {\v};
            \fi
        }
      
        \foreach \v [count=\k from 0] in {#3} {
            \pgfmathsetmacro\x{mod(\k,\rows)}
            \pgfmathsetmacro\y{int(\k/\rows)}
            \node[draw, minimum width=1cm, minimum height=1cm, align=center] at (\x,-\y) {\v};
        }

        \begin{scope}
            \clip (-0.5,0.5) rectangle (3.5,-3.5);
            \foreach \r in {#1} {\draw[rounded corners=1ex] \r;}
        \end{scope}

        \global\let\Vdim\undefined
    \end{tikzpicture}
}
\title{Digitaltechnik}
\institute{Karlsruher Institut für Technologie}

\author{\lqfullname}
\email{\lqemail @student.kit.edu}% chktex 1
% credit flo
\addtobeamertemplate{title page}{}{%
\scriptsize\color{gray}\begin{tikzpicture}[overlay, remember picture]
    \node[above left = 1em and 1em of current page.south east,align=center] (n0) {\LaTeX-Vorlage von\\\href{https://github.com/EagleoutIce}{Florian Sihler}};
    \node[above = 0cm of n0] (n1) {\fancyqr[size=1.5cm,tight]{https://eagleoutice.github.io/beamer-themes/\#colorful-dream}};
\end{tikzpicture}}% chktex 31
\usepackage{tikz}
\usetikzlibrary{fadings}

% cursed braces
\makeatletter
\newcommand{\cb@dx}[1]{%
    \ifnum`#1=`(  .0652em\else% chktex 37
    \ifnum`#1=`) -.0652em\else
    \ifnum`#1=`<  .0652em\else
    \ifnum`#1=`> -.0652em\else
    \ifnum`#1=`B  .0113em\else
    0em\fi\fi\fi\fi\fi
}
\newcommand{\cb@r}[1]{%
    \ifnum`#1=`< \langle\else
    \ifnum`#1=`> \rangle\else
    \ifnum`#1=`& \backslash\else
    \ifnum`#1=`= \forall\else% im not sure about this
    #1\fi\fi\fi\fi
}
\newcommand{\cb@clip}[4][]{\scoped{%
    \clip(-.3117em,#2) rectangle (.3117em,#3);%
    \node[#1] at (\cb@dx{#4},0em) {\(\cb@r{#4}\)};% chktex 1
}}
\newcommand{\cb@ht}[1][]{\cb@clip[#1]{0em}{.5669em}}
\newcommand{\cb@hb}[1][]{\cb@clip[#1]{-.5669em}{0em}}
\newcommand{\cb@tt}[1][]{\cb@clip[#1]{.1418em}{.5669em}}
\newcommand{\cb@tm}[1][]{\cb@clip[#1]{-.1418em}{.1418em}}
\newcommand{\cb@tb}[1][]{\cb@clip[#1]{-.5669em}{-.1418em}}
\newcommand{\cb@split}[4][]{%
    \ifx\relax#4\relax
        \cb@ht[#1]{#2}\cb@hb[#1]{#3}%
    \else
        \cb@tt[#1]{#2}\cb@tm[#1]{#3}\cb@tb[#1]{#4}%
    \fi
}
\newcommand{\cb}[2][]{{\tikz[baseline=-.8ex]{\cb@split[#1]#2\relax}}}
\newcommand{\cblogo}{\cb[inner sep=0pt,shading=axis,left color=orange,right color=violet,shading angle=0]{cB}}
\makeatother
% chktex-file 1 chktex-file 14 chktex-file 10 chktex-file 41 chktex-file 36 chktex-file 31

%% This is file `kvmacros.tex',
%% Version of January 7th, 2002
%% 
%% Copyright (C) 1998-2002 by Andreas W. Wieland, awwieland@gmx.de 
%%
%% This program can be redistributed and/or modified under the terms
%% of the LaTeX Project Public License distributed from CTAN
%% archives in directory macros/latex/base/lppl.txt; either
%% version 1 of the License, or (at your option) any later version,
%% with `The Program' referring to the software `kvmacros.tex' and its
%% accompanying documentation and `The Copyright Holder' referring to the
%% person Andreas W. Wieland.
%% 
%% 
%% 
%% IMPORTANT NOTICE: 
%% 
%% For error reports, comments or suggestions in case of UNCHANGED 
%% versions send mail to:
%% awwieland@gmx.de
%% 
%% Anyway, if you use my macros, send me an eMail, too. I would like to know
%% if they are useful to somebody out there.
%%
%% Please do not request updates from me directly. Distribution is 
%% done through Mail-Servers and TeX organizations. 
%% 
%% You are allowed to distribute this file under the condition that 
%% it is distributed together with `kvdoc.tex', `kvdoc.dvi' and `kvdoc.ps'. 
%% 
%% If you receive this file alone from someone, complain! 
%%
\typeout{}
\typeout{Macros for typesetting Karnaugh maps and Veitch charts}
\typeout{Version of January 7th, 2002}
\typeout{by Andreas W. Wieland, awwieland@gmx.de}
\typeout{}
%%
%% Change History:
%% August 04th, 1998: Original Version
%%
%% August 19th, 1998: Minor changes to make the lines within the maps look
%% better, also moved the variable identifiers from the right to the left and
%% added an identifier for the logic function; since then \kvmap has 5
%% instead of 4 parameters. I also included the documentation in an additional
%% Postscript-file.  
%%
%% September 23rd, 1998: Fixed the problem with double braces by introducing
%% the \compdummy and \ende macros in the \ifx-comparisons (see below): 
%% The parameters that are longer than one character do not have to be
%% enclosed in double braces any longer --- single braces are
%% sufficient. However, double braces will still be processed correctly.
%%
%% May 19th, 1999: Changed the documentation (which was months overdue) after
%% F. M. Brown pointed me to the incorrectness of the term
%% "Karnaugh-Veitch-map". Actually, the macros itself haven't
%% changed. Probably I should introduce macros for typesetting Veitch charts
%% as well (which would be a minor task). This version was actually never
%% released. 
%%
%% May 25th, 1999: Introduced macros for Veitch charts
%% (\veitchfoo) and renamed the macros that are solely related to
%% Karnaugh maps from \kvfoo to \karnaughfoo. Introduced aliases for those
%% macros that are user-callable to maintain compatibility with older versions:
%% \kvmap -> \karnaughmap. However, \kvindex, \kvnoindex and \kvunitlength
%% haven't changed their names as they are related to both Karnaugh maps and
%% Veitch charts. Changed the (plain stupid) algorithm that processes the
%% variable-length arguments so that it takes much less time to produce large
%% maps. Several macros so became redundant and are no longer included. Also I
%% introduced \kvindexsize and \kvcontentsize to adjust the font size of the
%% indices and the contents of a diagram. They default to \tiny and
%% \normalsize. 
%%
%% January 17th, 2001: Changed the eMail-address in the macro file and the
%% accompanying documentation and made minor changes to the documentation.
%%
%% January 7th, 2002: Exchanged the \xdef to \gdef in the definition of
%% \kvargumentstring and \kvgetonechar so that things like {\textbf 1} in the
%% argumentlist of \karnaughmap do not cause error messages any longer. 
%% Thanks to Michal Medveck{\'y} for pointing me to this error.
%%
%%
%%
%% We need a fixed dimension for a single field in a Karnaugh map:
%%
\newdimen\kvunitlength
\kvunitlength=8mm
%%
%% We need a default font size for the indices:
%%
\newcommand{\kvindexsize}{\tiny}
%%
%% And we need a default size for the contents and variable identifiers:
%%
\newcommand{\kvcontentsize}{\normalsize}
%%
%% First, we have to introduce some counters:
%%
%% \kvrecursiondepth is used to control the recursion of the
%% \karnaughmakemap- and \veitchmakechart-macro. 
%%
\newcount\kvrecursiondepth
%%
%% The \kvindexcounter is needed for the indices in the fields of the
%% diagrams. 
%%
\newcount\kvindexcounter
%% 
%% \kvxsize and \kvysize store the dimensions of an entire diagram.
%%
\newcount\kvxsize
\newcount\kvysize
%%
%% Some counters are necessary to compute the marks for the variable
%% identifiers:
%%
\newcount\kvvarno
\newcount\kvxvarno
\newcount\kvyvarno
\newcount\kvmarkstart
\newcount\kvmarklength
\newcount\kvmarknum
\newcount\kvmarkmove
%%
%% And we need a savebox to store the variable marks:
%%
\newsavebox\kvsavebox
%%
%% Single fields in a diagram should be indexed, which makes the map easier to
%% use. This is the default. If you don't want indices, simply call
%% \kvnoindex. If you want indices back, call \kvindex:
%%
\def\kvnoindex{%
\def\kvcurrentindex{}}
%%
\def\kvindex{%
\def\kvcurrentindex{%
\the\kvindexcounter\global\advance\kvindexcounter by 1}%
}
%%
\kvindex
%%
%% We need a macro that computes the powers of two:
%%
\def\kvpoweroftwo#1#2{% Computes #1=2^#2, both of which have to be counters
{\ifnum#2>0 
\global\multiply#1 by 2 
\advance#2 by -1
\kvpoweroftwo{#1}{#2}
\fi}}
%%
%% The macros \kvargumentstring, \kvgetchar and \kvgetonechar are needed to
%% process the variabe-length parameters in \karnaughmap and \veitchchart:
%%
\def\kvargumentstring#1{\gdef\kvdummystring{#1{}\noexpand\end}}
%%
\def\kvgetchar{\expandafter\kvgetonechar\kvdummystring}
%%
\def\kvgetonechar#1#2\end{{#1}\gdef\kvdummystring{#2\noexpand\end}}%
%%
%% The Macro \karnaughmakemap calls itself recursively until the parameter #1
%% equals 1, whereupon it returns a single token from the list of arguments,
%% enclosed in a \makebox plus the index (if enabled) in a smaller \makebox:
%%
\def\karnaughmakemap#1#2{{%
\kvrecursiondepth=\number#1
\ifnum\kvrecursiondepth>1
\divide\kvrecursiondepth by 2
\unitlength=\kvunitlength
\multiply\unitlength by \kvrecursiondepth
%%
\ifcase#2
%%
%% The parameter #2 of \karnaughmakemap is needed because the inner Karnaugh
%% maps need to be mirrored. This is achieved by the following case-statement,
%% which orders the inner Karnaugh maps properly:
%%
%% Case 0: top-left Karnaugh map
\begin{picture}(2,2)%
\put(0,1){\karnaughmakemap{\kvrecursiondepth}{0}}%
\put(1,1){\karnaughmakemap{\kvrecursiondepth}{1}}%
\put(0,0){\karnaughmakemap{\kvrecursiondepth}{2}}%
\put(1,0){\karnaughmakemap{\kvrecursiondepth}{3}}%
\end{picture}%
\or
%% Case 1: top-right Karnaugh map
\begin{picture}(2,2)%
\put(1,1){\karnaughmakemap{\kvrecursiondepth}{1}}%
\put(0,1){\karnaughmakemap{\kvrecursiondepth}{0}}%
\put(1,0){\karnaughmakemap{\kvrecursiondepth}{3}}%
\put(0,0){\karnaughmakemap{\kvrecursiondepth}{2}}%
\end{picture}%
\or
%% Case 2: bottom-left Karnaugh map
\begin{picture}(2,2)%
\put(0,0){\karnaughmakemap{\kvrecursiondepth}{2}}%
\put(1,0){\karnaughmakemap{\kvrecursiondepth}{3}}%
\put(0,1){\karnaughmakemap{\kvrecursiondepth}{0}}%
\put(1,1){\karnaughmakemap{\kvrecursiondepth}{1}}%
\end{picture}%
\or
%% Case 3: bottom-right Karnaugh map
\begin{picture}(2,2)%
\put(1,0){\karnaughmakemap{\kvrecursiondepth}{3}}%
\put(0,0){\karnaughmakemap{\kvrecursiondepth}{2}}%
\put(1,1){\karnaughmakemap{\kvrecursiondepth}{1}}%
\put(0,1){\karnaughmakemap{\kvrecursiondepth}{0}}%
\end{picture}%
\fi
\else
\unitlength=\kvunitlength
\begin{picture}(1,1)
\put(0,0){\makebox(1,1){\kvcontentsize\kvgetchar}}%
\put(0.05,0.05){\makebox(0.9,0.9)[bl]{\kvindexsize\kvcurrentindex}}%
\end{picture}
\fi}}%
%%
%% \karnaughmaketopmark typesets the variable marks of a Karnaugh map that are
%% located on top of the diagram:
%%
\def\karnaughmaketopmark{%
  \unitlength\kvunitlength
  \begin{picture}(\kvxsize,1)
    \kvmarkstart=1
    \kvpoweroftwo{\kvmarkstart}{\kvxvarno} % \kvymarkstart is the start 
                                % position for the \multiput
    \kvmarklength=\kvmarkstart 
    \multiply\kvmarklength by 2 % \kvmarklength is the length of a mark
    \kvmarkmove=\kvmarkstart 
    \multiply\kvmarkmove by 4 % This is the move distance for the \multiput.
    \kvmarknum=\kvxsize
    \divide\kvmarknum by \kvmarkmove % This is the number of repetitions for
                                % the \multiput.  
    %The highest-order variable mark needs a special treatment:
    \ifnum\kvmarknum=0\kvmarknum=1\divide\kvmarklength by 2\fi 
    \savebox\kvsavebox(\kvmarklength,1){% 
      \begin{picture}(\kvmarklength,1)
      \put(0,0.3){\makebox(\kvmarklength,0.7){\kvcontentsize\kvgetchar}} 
      \put(0,0.1){\line(0,1){0.4}} 
      \put(\kvmarklength,0.1){\line(0,1){0.4}}
      \put(0,0.3){\line(1,0){\kvmarklength}} 
      \end{picture}}
    \multiput(\kvmarkstart,0)(\kvmarkmove,0){\kvmarknum}{\usebox\kvsavebox}
  \end{picture}
}
%%
%% \karnaughmakeleftmark typesets the variable marks of a Karnaugh map that are
%% located on the left of the diagram:
%%
\def\karnaughmakeleftmark{%
  \unitlength\kvunitlength
  \begin{picture}(-1,\kvysize)(0,-\kvysize)
    \kvmarkstart=1
    \kvpoweroftwo{\kvmarkstart}{\kvyvarno} % \kvmarkstart is the start 
                                % position for the \multiput
    \kvmarklength=\kvmarkstart 
    \multiply\kvmarklength by 2 % \kvmarklength is the length of a mark
    \kvmarkmove=\kvmarkstart 
    \multiply\kvmarkmove by 4 % This now is the move distance for the
                              % \multiput. 
    \kvmarknum=\kvysize
    \divide\kvmarknum by \kvmarkmove % This now is the number of 
                                % repetitions for the \multiput.  
    %The highest-order variable mark needs a special treatment:
    \ifnum\kvmarknum=0\kvmarknum=1\divide\kvmarklength by 2\fi 
    \advance\kvmarkstart by \kvmarklength
    \savebox\kvsavebox(1,\kvmarklength){%
      \begin{picture}(1,\kvmarklength)
        \put(-0.3,0){\makebox(-0.7,\kvmarklength){\kvcontentsize\kvgetchar}}
        \put(-0.1,0){\line(-1,0){0.4}}
        \put(-0.1,\kvmarklength){\line(-1,0){0.4}}
        \put(-0.3,0){\line(0,1){\kvmarklength}}
      \end{picture}}
    \multiput(0,-\kvmarkstart)(0,-\kvmarkmove){\kvmarknum}{\usebox\kvsavebox}
  \end{picture}
}
%% \karnaughmakemarks calls \karnaughmaketopmark or \karnaughmakeleftmark
%% depending on whether \kvvarno is odd or even.
%%
\def\karnaughmakemarks{%
\ifnum\kvvarno>0
  \let\next=\karnaughmakemarks
  \ifodd\kvvarno % We have to make a mark at the top
    \advance\kvxvarno by -1
    \put(0,\kvxvarno){\karnaughmaketopmark}
  \else % We have to make a mark at the left
    \advance\kvyvarno by -1
    \put(-\kvyvarno,-\kvysize){\karnaughmakeleftmark}
  \fi
  \advance\kvvarno by -1
\else
  \let\next=\relax
\fi
\next
}
%%
%% \karnaughmap is the macro that a user calls if he wants to draw a
%% Karnaugh map:  
%%
\def\karnaughmap#1#2#3#4#5{%
%%
%% #1 is the number of variables in the Karnaugh map
%% #2 is the identifier of the function
%% #3 is the list of identifiers of those variables 
%% #4 is the list of tokens that have to be written into the map
%% #5 is something that you want to draw inside the Karnaugh map
%%
\kvvarno=#1  % \kvvarno is the total number of variables 
\kvyvarno=#1 % \kvyvarno is the number of variable marks at the left 
\divide\kvyvarno by 2
\kvxvarno=#1 % \kvxvarno is the number of variable marks on top 
\advance\kvxvarno by -\kvyvarno
\kvxsize=1
\kvpoweroftwo{\kvxsize}{\kvxvarno}
\kvysize=1
\kvpoweroftwo{\kvysize}{\kvyvarno}
\advance\kvxsize by \kvyvarno
\advance\kvxvarno by \kvysize
\unitlength\kvunitlength
\begin{picture}(\kvxsize,\kvxvarno)(-\kvyvarno,-\kvysize)
\advance\kvxsize by -\kvyvarno
\advance\kvxvarno by -\kvysize
\put(0,-\kvysize){%
\begin{picture}(\kvxsize,\kvysize)
\multiput(0,0)(0,1){\kvysize}{\line(1,0){\kvxsize}}
\multiput(0,0)(1,0){\kvxsize}{\line(0,1){\kvysize}}
\put(0,\kvysize){\line(1,0){\kvxsize}}
\put(\kvxsize,0){\line(0,1){\kvysize}}
#5
\end{picture}}
\put(-\kvyvarno,0){\makebox(\kvyvarno,\kvxvarno){#2}}
\kvindexcounter=0
\kvargumentstring{#4}
\put(0,-\kvysize){\karnaughmakemap{\kvysize}{0}}
\ifodd\kvvarno
{\divide\kvxsize by 2
\put(\kvxsize,-\kvysize){\karnaughmakemap{\kvysize}{1}}}
\fi
\kvargumentstring{#3}
\karnaughmakemarks
%%
\end{picture}
}%
%%
%% The definition of \kvmap is necessary to maintain compatibility with older
%% versions of this macro package:
%%
\def\kvmap{\karnaughmap}%
%%
%%
%% The Macro \veitchmakechart calls itself recursively until the parameter #1
%% equals 1, whereupon it returns a single token from the list of arguments,
%% enclosed in a \makebox plus the index (if enabled) in a smaller \makebox:
%%
\def\veitchmakechart#1{{%
\kvrecursiondepth=\number#1
\ifnum\kvrecursiondepth>1
\divide\kvrecursiondepth by 2
\unitlength=\kvunitlength
\multiply\unitlength by \kvrecursiondepth
%%
\begin{picture}(2,2)%
\put(0,1){\veitchmakechart{\kvrecursiondepth}}%
\put(1,1){\veitchmakechart{\kvrecursiondepth}}%
\put(0,0){\veitchmakechart{\kvrecursiondepth}}%
\put(1,0){\veitchmakechart{\kvrecursiondepth}}%
\end{picture}%
\else
\unitlength=\kvunitlength
\begin{picture}(1,1)
\put(0,0){\makebox(1,1){\kvcontentsize\kvgetchar}}%
\put(0.05,0.05){\makebox(0.9,0.9)[bl]{\kvindexsize\kvcurrentindex}}%
\end{picture}
\fi}}%
%%
\def\veitchmaketopmark{%
  \unitlength\kvunitlength
  \begin{picture}(\kvxsize,1)
    \kvmarkstart=1
    \kvpoweroftwo{\kvmarkstart}{\kvxvarno} % \kvymarkstart is the start 
                                % position for the \multiput
    \kvmarklength=\kvmarkstart % \kvmarklength is the length of a mark
    \kvmarkmove=\kvmarkstart 
    \multiply\kvmarkmove by 2 % This is the move distance for the \multiput.
    \kvmarknum=\kvxsize
    \divide\kvmarknum by \kvmarkmove % This is the number of repetitions for
                                     % the \multiput.  
    \savebox\kvsavebox(\kvmarklength,1){% 
      \begin{picture}(\kvmarklength,1)
      \put(0,0.3){\makebox(\kvmarklength,0.7){\kvcontentsize\kvgetchar}} 
      \put(0,0.1){\line(0,1){0.4}} 
      \put(\kvmarklength,0.1){\line(0,1){0.4}}
      \put(0,0.3){\line(1,0){\kvmarklength}} 
      \end{picture}}
    \multiput(\kvmarkstart,0)(\kvmarkmove,0){\kvmarknum}{\usebox\kvsavebox}
  \end{picture}
}
%%
\def\veitchmakeleftmark{%
  \unitlength\kvunitlength
  \begin{picture}(-1,\kvysize)(0,-\kvysize)
    \kvmarkstart=1
    \kvpoweroftwo{\kvmarkstart}{\kvyvarno} % \kvmarkstart is the start 
                                % position for the \multiput
    \kvmarklength=\kvmarkstart 
    \kvmarkmove=\kvmarkstart 
    \multiply\kvmarkmove by 2 % This now is the move distance for the \multiput. 
    \kvmarknum=\kvysize
    \divide\kvmarknum by \kvmarkmove % This now is the number of 
                                     % repetitions for the \multiput.  
    \advance\kvmarkstart by \kvmarklength
    \savebox\kvsavebox(1,\kvmarklength){%
      \begin{picture}(1,\kvmarklength)
        \put(-0.3,0){\makebox(-0.7,\kvmarklength){\kvcontentsize\kvgetchar}}
        \put(-0.1,0){\line(-1,0){0.4}}
        \put(-0.1,\kvmarklength){\line(-1,0){0.4}}
        \put(-0.3,0){\line(0,1){\kvmarklength}}
      \end{picture}}
    \multiput(0,-\kvmarkstart)(0,-\kvmarkmove){\kvmarknum}{\usebox\kvsavebox}
  \end{picture}
}
%%
\def\veitchmakemarks{%
\ifnum\kvvarno>0
  \let\next=\veitchmakemarks
  \ifodd\kvvarno % We have to make a mark at the top
    \advance\kvxvarno by -1
    \put(0,\kvxvarno){\veitchmaketopmark}
  \else % We have to make a mark at the left
    \advance\kvyvarno by -1
    \put(-\kvyvarno,-\kvysize){\veitchmakeleftmark}
  \fi
  \advance\kvvarno by -1
\else
  \let\next=\relax
\fi
\next
}
%%
\def\veitchchart#1#2#3#4#5{%
%%
%% #1 is the number of variables in the Veitch chart
%% #2 is the identifier of the function
%% #3 is the list of identifiers of those variables 
%% #4 is the list of tokens that have to be written into the chart
%% #5 is something that you want to draw inside the Veitch chart
%%
\kvvarno=#1  % \kvvarno is the total number of variables 
\kvyvarno=#1 % \kvyvarno is the number of variable marks at the left 
\divide\kvyvarno by 2
\kvxvarno=#1 % \kvxvarno is the number of variable marks on top 
\advance\kvxvarno by -\kvyvarno
\kvxsize=1
\kvpoweroftwo{\kvxsize}{\kvxvarno}
\kvysize=1
\kvpoweroftwo{\kvysize}{\kvyvarno}
\advance\kvxsize by \kvyvarno
\advance\kvxvarno by \kvysize
\unitlength\kvunitlength
\begin{picture}(\kvxsize,\kvxvarno)(-\kvyvarno,-\kvysize)
\advance\kvxsize by -\kvyvarno
\advance\kvxvarno by -\kvysize
\put(0,-\kvysize){%
\begin{picture}(\kvxsize,\kvysize)
\multiput(0,0)(0,1){\kvysize}{\line(1,0){\kvxsize}}
\multiput(0,0)(1,0){\kvxsize}{\line(0,1){\kvysize}}
\put(0,\kvysize){\line(1,0){\kvxsize}}
\put(\kvxsize,0){\line(0,1){\kvysize}}
#5
\end{picture}}
\put(-\kvyvarno,0){\makebox(\kvyvarno,\kvxvarno){#2}}
\kvindexcounter=0
\kvargumentstring{#4}
\put(0,-\kvysize){\veitchmakechart{\kvysize}}
\ifodd\kvvarno
{\divide\kvxsize by 2
\put(\kvxsize,-\kvysize){\veitchmakechart{\kvysize}}}
\fi
\kvargumentstring{#3}
\veitchmakemarks
%%
\end{picture}
}%
%%

%%% Local Variables: 
%%% mode: latex
%%% TeX-master: t
%%% End: 
% Vorlage von VL

\usepackage{circuitikz}% chktex-file 1 chktex-file 8 chktex-file 26
\usetikzlibrary{overlay-beamer-styles}

%%%%%%%%%%%%%%%%%%%%%%%%%%%%%%%%%%%%%%%%%%%%%%%%%%%%%%%%%%%%%%%%

\subtitle{Tutorium 8}
\date{07. Juli 2026}

\titleimagebox{1}{7cm}{\begin{tabular}[t]{c|ccccccc}%chktex 44
      & Q & U & I & N & E &   &   \\
    \hline \hline%chktex 44
    M &   &   &   &   &   &   & \cb{"\{U} \\% chktex 18
    C & C & L & U & S & K & E & Y \\
  \end{tabular}}

\outroimagetext{Wann kann GPT endlich mir die Diagramme zeichnen?}

%%%%%%%%%%%%%%%%%%%%%%%%%%%%%%%%%%%%%%%%%%%%%%%%%%%%%%%%%%%%%%%%

\begin{document}

\section[Böse Böse Böse]{Hinweis zu KI}
\begin{frame}{KV-Diagramme sehen \textbf{so} aus:}\onslide<+->
    \begin{columns}[T]
    \begin{column}{0.5\textwidth}
        \karnaughmap{3}{\(f(z, y, x)\):}{{\(z\)}{\(y\)}{\(x\)}}{01010010}{}
    \end{column}
    \begin{column}{0.5\textwidth}
        \onslide<+->\textbf{NICHT}\vspace*{2\baselineskip}

        \begin{tabular}{c|cccc}% chktex 44
            & 00 & 01 & 10 & 11 \\
            \hline%chktex 44
            0 & 0 & 1 & 0 & 0 \\
            1 & 0 & 1 & 1 & 0 \\
        \end{tabular}

        {\tiny Danke GPT}
    \end{column}
    \end{columns}\vspace*{2\baselineskip}

    \onslide<+->Ab jetzt: \alert{0 Punkte}!
\end{frame}

\section{Korrektur}

\subsection{Aufgabe 1}
\begin{frame}{Aufgabe 1}\onslide<+->
    Zur Übung an der Tafel.\onslide<+->~Und damit ich weniger vorbereiten muss.\onslide<+->
    \begin{exercise}
        \(h(d,c,b,a)=d\overline c\overline ba\lor d\overline cba\lor dc\overline b\lor\overline dcb\lor dcb\)
    \end{exercise}
    \onslide<+->
    \begin{solve}[DMF]
        \(cb \lor da \lor dc\)
    \end{solve}
\end{frame}

\subsection{Aufgabe 2}
\begin{frame}{Dürfen wir das?}\onslide<+->
    \begin{exercise}[%
        \only<2,3>{\(A\) streichen?}\only<3>{~\faAngleRight~Ja, da Kern}%
        \only<4,5>{\(B\) streichen?}\only<5>{~\faAngleRight~Ja, da überdeckt}%
        \only<6,7>{\(F\) streichen?}\only<7>{~\faAngleRight~Ja, da Kern}%
        \only<8,9>{\(6\) streichen?}\only<9>{~\faAngleRight~Nein, keine Spaltendominanz}%
        \only<10,11>{\(7\) streichen?}\only<11>{~\faAngleRight~Ja, da von \(6\) dominiert}%
        ]
        \begin{tabular}[t]{c|ccccccccc}%chktex 44
            Prim- & \multicolumn{9}{c}{Minterme der Funktion \(h\)} \\
            implikanten & 0 & 1 & 2 & 3 & 4 & \alert<9,11>{6} & \alert<11>{7} & 13 & 15 \\ \hline \hline%chktex 44
            \(A\) & & & \(\times \) & \(\times \) & \(\times \) & \(\alert<3,5,9,11>{\times} \) & \(\alert<5,11>{\times} \) \\
            \(B\) & & & \(\times \) & \(\times \) \\
            \(C\) & \(\times \) & \(\times \) & & & \(\times \) \\
            \(D\) & \(\times \) & \(\times \) \\
            \(E\) & \(\times \) & & \(\times \) & & \(\times \) & & \(\alert<11>{\times} \) \\
            \(F\) & & & & & & & & \(\alert<7>{\times} \) & \(\alert<7>{\times} \) \\
          \end{tabular}
    \end{exercise}
    \begin{solve}[1.]
        \only<3,7>{Markieren nicht vergessen!}%
        \only<5>{Nur wenn Kosten von \(A\) geringer als von \(B\)!}%
        \only<9>{Wenn ein Kernprimimplikant (wie \(A\)) gestrichen wird, dann natürlich schon!}%
    \end{solve}
\end{frame}
\resetsolve%

\begin{frame}{Und vereinfachen}\onslide<+->
    \begin{solve}[2.: Überdeckungsfunktion]\onslide<+->
        \(A, F\) Kern, Übrig: \(C, D, E, 0, 1\)

        \onslide<+->\((w_C\lor w_D\lor w_E)\land(w_C\lor w_D)\)
    \end{solve}
    \resetsolve\onslide<+->
    \begin{solve}[3.: DMF]
        \onslide<+->\(A\lor C\lor F\) oder \(A\lor D\lor F\)
    \end{solve}
    \resetsolve\onslide<+->
    \begin{solve}[4.: Vollständig?]
        \onslide<+->Nein, da \(C\) nur 3 Minterme überdeckt (keine Zweierpotenz)
    \end{solve}
\end{frame}

\subsection{Aufgabe 3}
\begin{frame}{Ich packe in mein Bündel\dots}\onslide<+->
    \begin{exercise}
        \(f = \overline a \, \overline c \, \overline d \lor a \, \overline b c d \lor \overline a \, \overline c d\)

        \(g = a c \, \overline b \lor \overline a b \, \overline c \, \overline d\)

        \(h = \overline a c \lor \overline a d \lor \overline b c \, \overline d \lor \overline a \, \overline b \, \overline d \, \overline c\)
    \end{exercise}\onslide<+->
    \begin{solve}
    \only<.-+>{\resizebox{!}{0.2\textheight}{
    \begin{minipage}{5cm}
        \karnaughmap{4}{\(f\):}{{\(d\)}{\(c\)}{\(b\)}{\(a\)}}{1010000010100100}{\only<.>{
            \put(0.5, 2){\textcolor{blue}{\oval(0.9, 3.9)}}
            \put(2.5, 0.5){\textcolor{red}{\oval(0.9, 0.9)}}
        }}
    \end{minipage}
    \begin{minipage}{5cm}
        \karnaughmap{4}{\(g\):}{{\(d\)}{\(c\)}{\(b\)}{\(a\)}}{0010010000000100}{\only<.>{
            \put(0.5, 2.5){\textcolor{blue}{\oval(0.9, 0.9)}}
            \put(2.5, 0){\textcolor{red}{\oval(0.9, 1.9)[t]}}
            \put(2.5, 4){\textcolor{red}{\oval(0.9, 1.9)[b]}}
        }}
    \end{minipage}
    \begin{minipage}{5cm}
        \karnaughmap{4}{\(h\):}{{\(d\)}{\(c\)}{\(b\)}{\(a\)}}{1000111010101010}{\only<.>{
            \put(3.5, 2){\textcolor{blue}{\oval(0.9, 3.9)}}
            \put(0, 1){\textcolor{red}{\oval(1.9, 1.9)[r]}}
            \put(4, 1){\textcolor{red}{\oval(1.9, 1.9)[l]}}
            \put(4, 1){\textcolor{red}{\oval(1.9, 1.9)[l]}}
            \put(3, 3.5){\textcolor{orange}{\oval(1.8, 0.8)}}
            \put(0, 0){\textcolor{black}{\oval(1.9, 1.9)[rt]}}
            \put(4, 0){\textcolor{black}{\oval(1.9, 1.9)[lt]}}
            \put(0, 4){\textcolor{black}{\oval(1.9, 1.9)[rb]}}
            \put(4, 4){\textcolor{black}{\oval(1.9, 1.9)[lb]}}
        }}
    \end{minipage}
    }}
    \only<+>{
    \(f = \overline a \, \overline c\lor a \, \overline b c d\)

    \(g = a c \, \overline b \lor \overline a b \, \overline c \, \overline d\) alle Kern

    \(h = \overline a \, \overline c \lor \overline a d \lor \overline b c \, \overline d \lor \overline a \, \overline b\)
    
    Insgesamt \textbf{8} Terme, aber es gibt Koppelterme!}
    \only<+>{\resizebox{!}{0.2\textheight}{
    \begin{minipage}{5cm}
        \karnaughmap{4}{\(f\):}{{\(d\)}{\(c\)}{\(b\)}{\(a\)}}{1010000010100100}{
            \put(0.5, 4){\textcolor{blue}{\oval(0.9, 1.9)[b]}}
            \put(0.5, 0){\textcolor{blue}{\oval(0.9, 1.9)[t]}}
            \put(0.5, 2.5){\textcolor{red}{\oval(0.9, 0.9)}}
            \put(0.5, 1){\textcolor{black}{\oval(0.9, 1.9)}}
            \put(2.5, 0.5){\textcolor{green}{\oval(0.9, 0.9)}}
        }
    \end{minipage}
    \begin{minipage}{5cm}
        \karnaughmap{4}{\(g\):}{{\(d\)}{\(c\)}{\(b\)}{\(a\)}}{0010010000000100}{
            \put(0.5, 2.5){\textcolor{red}{\oval(0.9, 0.9)}}
            \put(2.5, 0.5){\textcolor{green}{\oval(0.9, 0.9)}}
            \put(2.5, 3.5){\textcolor{orange}{\oval(0.9, 0.9)}}
        }
    \end{minipage}
    \begin{minipage}{5cm}
        \karnaughmap{4}{\(h\):}{{\(d\)}{\(c\)}{\(b\)}{\(a\)}}{1000111010101010}{
            \put(0.5, 4){\textcolor{blue}{\oval(0.9, 1.9)[b]}}
            \put(0.5, 0){\textcolor{blue}{\oval(0.9, 1.9)[t]}}
            \put(0.5, 1){\textcolor{black}{\oval(0.9, 1.9)}}
            \put(2.5, 3.5){\textcolor{orange}{\oval(0.9, 0.9)}}
            \put(3.5, 2){\textcolor{magenta}{\oval(0.9, 3.9)}}
        }
    \end{minipage}
    }}
\end{solve}
\end{frame}

\subsection{Aufgabe 4}
\begin{frame}{Teile den Kuchen}\onslide<+->
    \begin{columns}[T]
    \begin{column}{0.5\textwidth}
    \begin{solve}[Funktionstabelle]
    \begin{tabular}[t]{ccc|cccc}%chktex 44
        \(c\) & \(b\) & \(a\) & \(t2\) & \(t3\) & \(t5\) & \(t7\) \\
        \hline%chktex 44
        0 & 0 & 0 & \only<+->{- & - & - & -} \\
        0 & 0 & 1 & \only<+->{0 & 0 & 0 & 0} \\
        0 & 1 & 0 & \only<+->{1 & 0 & 0 & 0} \\
        0 & 1 & 1 & \only<+->{0 & 1 & 0 & 0} \\
        1 & 0 & 0 & \only<+->{1 & 0 & 0 & 0} \\
        1 & 0 & 1 & \only<+->{0 & 0 & 1 & 0} \\
        1 & 1 & 0 & \only<+->{1 & 1 & 0 & 0} \\
        1 & 1 & 1 & \only<+->{0 & 0 & 0 & 1}
    \end{tabular}
    \end{solve}
    \end{column}
    \resetsolve\onslide<+->
    \begin{column}{0.5\textwidth}
    \begin{solve}[Schaltbild]
    \resizebox{\textwidth}{!}{
    \begin{tikzpicture}[font=\small]

        %-------------------------------------------------
        % Decoderblock
        %-------------------------------------------------
        \draw (0,0) rectangle (3.3,6);
        
        \node at (1.65,5.7) {\Large DX};
        
        % Eingänge
        \draw (-0.6,4.8) -- (0,4.8);
        \draw (-0.6,4.0) -- (0,4.0);
        \draw (-0.6,3.2) -- (0,3.2);
        
        \node[left] at (-0.6,4.8) {$a$};
        \node[left] at (-0.6,4.0) {$b$};
        \node[left] at (-0.6,3.2) {$c$};
        
        \node[right] at (0.15,4.8) {$0$};
        \node[right] at (0.15,4.0) {$1$};
        \node[right] at (0.15,3.2) {$2$};
        
        % Klammer und Bezeichnung
        \node at (0.85,4.0) {\Huge $\}$};
        \node at (1.55,4.0) {$G^{0}_{7}$};
        
        % Enable
        \draw (-0.6,1.0) -- (0,1.0);
        \node[left] at (-0.6,1.0) {$1$};
        \node at (0.6,1.0) {EN};
        
        % Ausgänge 0..7
        \foreach \n/\y in {
        0/5.0,
        1/4.2,
        2/3.4,
        3/2.6,
        4/2.0,
        5/1.4,
        6/0.8,
        7/0.2}
        {
            \draw (3.3,\y) -- (3.6,\y);
            \node[left] at (3.28,\y) {\n};
        }
        
        %-------------------------------------------------
        % OR-Block t2  (>=1)
        %-------------------------------------------------
        \draw (5.1,3.8) rectangle (6.2,4.9);
        \node at (5.65,4.35) {$\geq 1$};
        \node[right] at (6.55,4.35) {$t2$};
        \draw (6.2,4.35) -- (6.5,4.35);
        
        %-------------------------------------------------
        % OR-Block t3  (>=1)
        %-------------------------------------------------
        \draw (5.1,2.0) rectangle (6.2,3.1);
        \node at (5.65,2.55) {$\geq 1$};
        \node[right] at (6.55,2.55) {$t3$};
        \draw (6.2,2.55) -- (6.5,2.55);
        
        %-------------------------------------------------
        % Verdrahtung
        %-------------------------------------------------
        
        % Ausgang 2 -> oberes OR
        \draw (3.6,3.4) -- (4.1,3.4) -- (4.1,4.7) -- (5.1,4.7);
        
        % Ausgang 4 -> oberes OR
        \draw (3.6,2.0) -- (4.3,2.0) -- (4.3,4.4) -- (5.1,4.4);

        % Ausgang 6 -> oberes OR
        \draw (3.6,0.8) -- (4.5,0.8) -- (4.5,4.1) -- (5.1,4.1);
        \fill (4.5,0.8) circle (1.5pt);
        
        % Ausgang 3 -> unteres OR
        \draw (3.6,2.6) -- (4.7,2.6) -- (4.7,2.8) -- (5.1,2.8);
        
        % Ausgang 6 -> unteres OR
        \draw (3.6,0.8) -- (4.9,0.8) -- (4.9,2.3) -- (5.1,2.3);
        
        % Ausgang 5 -> t5
        \draw (3.6,1.4) -- (6.5,1.4);
        \node[right] at (6.55,1.4) {$t5$};
        
        % Ausgang 7 -> t7
        \draw (3.6,0.2) -- (6.5,0.2);
        \node[right] at (6.55,0.2) {$t7$};
        
    \end{tikzpicture}
    }
    \end{solve}
    \end{column}
    \end{columns}
\end{frame}

\end{document}% chktex 17